\section{The combined pipeline}\label{sec:stack}

The families above improve different stages, so the best pieces can be chained. The recommended single-frame
pipeline (\algo{stack\_dp\_rd\_gate}) is shown in Figure~\ref{fig:pipeline}; its core is twenty lines of glue
(Listing~\ref{lst:stack}).

\begin{figure}[H]
\centering
\begin{tikzpicture}[node distance=5mm, every node/.style={font=\small},
  blk/.style={draw=accent, fill=boxbg, rounded corners=2pt, minimum height=11mm, text width=24mm, align=center},
  io/.style={draw=black!50, rounded corners=2pt, minimum height=11mm, text width=17mm, align=center},
  arr/.style={-{Stealth[length=2mm]}, thick, black!70}]
\node[io] (img) {image\\$320\times240$};
\node[blk, right=of img] (dp) {\textbf{1. Line}\\Viterbi tracking};
\node[blk, right=of dp] (cal) {\textbf{2. Calibrate}\\pixel $\to$ mm};
\node[blk, right=of cal] (rd) {\textbf{3. Depth}\\RANSAC surface, DBSCAN floor};
\node[blk, right=of rd] (gate) {\textbf{4. Gates}\\cap, frame gate, risk gate};
\node[io, right=of gate] (out) {depths, PASS / FAIL or \emph{retake}};
\draw[arr] (img)--(dp); \draw[arr] (dp)--(cal); \draw[arr] (cal)--(rd); \draw[arr] (rd)--(gate); \draw[arr] (gate)--(out);
\end{tikzpicture}
\caption{The recommended single-frame pipeline. Stages 1--3 are classical geometry; only the two gates in stage 4
are learned (logistic regression).}
\label{fig:pipeline}
\end{figure}

\begin{lstlisting}[caption={The composition (\texttt{algos/stackcommon.py}, shortened).},label={lst:stack}]
def measure(gray, calib=None, off=None, line_fn=None, cap=False, lb=False, frame_gate=False, risk_model=None):
    line = line_fn(gray, off)                                          # 1. dp_dense.extract_line
    out = robust_depth.stage(line["u"], line["v"], calib, **RD)        # 2. rd_final: floor=dbscan, detect=shape
    out["line"] = line
    if cap:                                                            # 3. deeper than 25 mm is not a groove
        apply_cap(out, gray.shape[0])
    out["verdict"] = verdict_rule(out["grooves"]) if calib is not None else None   # 4. PASS / FAIL / RETAKE
    usable = 1.0 if out["grooves"] else 0.0                            # 5. no groove found = retake
    if frame_gate and usable >= 0.5:                                   #    is there a laser line on a tyre at all?
        out["p_frame"], _ = _prob("frame_lr", gate_features.features(gray)[0])
        usable = min(usable, out["p_frame"])
    if risk_model and calib is not None and usable >= 0.5:             #    is one of these depths probably wrong?
        out["p_bad"], t = _prob(risk_model, stack_features(gray, line, out))
        if out["p_bad"] > t:
            usable = 0.4
    out["usable"] = usable                                             # < 0.5 -> "retake the photo"
    return out
\end{lstlisting}

\subsection{Ablation: what each piece is worth}

Table~\ref{tab:ablation} adds one piece at a time. Rows 1a and 1b each change one thing against the baseline; rows
3--6 build on row 2; row 7 swaps the line extractor of row 2 for laser on/off subtraction.

\begin{table}[H]
\centering\footnotesize\setlength{\tabcolsep}{3.4pt}
\caption{Ablation on \texttt{stress} (top) and on other holder units, \texttt{shifted} (bottom). Depth columns in
mm. Bold = best in column.}\label{tab:ablation}
\begin{tabular}{@{}lrrrrrrrrr@{}}
\toprule
Step & Cov. & MAE & Median & P95 & Gross & Miss & False & Verdict & F.pass \\
\midrule
0\quad Baseline & 1.00 & 1.516 & 0.535 & 6.678 & 0.364 & 0.099 & 0.210 & 0.915 & 0.015 \\
1a\quad DP line only & 1.00 & 0.544 & 0.086 & 0.591 & 0.033 & 0.018 & 0.115 & 0.955 & \textbf{0.000} \\
1b\quad Robust depth only & 1.00 & 0.464 & 0.093 & 2.225 & 0.121 & 0.113 & 0.145 & 0.940 & 0.010 \\
2\quad DP line + robust depth & 1.00 & 0.325 & 0.088 & 0.540 & 0.036 & 0.015 & 0.065 & 0.955 & \textbf{0.000} \\
3\quad + 25 mm cap & 1.00 & 0.325 & 0.088 & 0.540 & 0.036 & 0.015 & 0.065 & 0.955 & \textbf{0.000} \\
4\quad + lower-bound rule & 0.99 & 0.325 & 0.088 & 0.540 & 0.036 & 0.006 & 0.065 & 0.960 & \textbf{0.000} \\
5\quad + risk gate & 0.91 & 0.246 & 0.088 & 0.490 & 0.026 & 0.007 & \textbf{0.016} & \textbf{0.995} & \textbf{0.000} \\
6\quad + frame gate \textbf{(recommended)} & 0.89 & 0.249 & 0.088 & 0.497 & 0.027 & \textbf{0.003} & \textbf{0.017} & \textbf{0.994} & \textbf{0.000} \\
7\quad Row 2 with on/off line & 1.00 & \textbf{0.142} & \textbf{0.084} & \textbf{0.339} & \textbf{0.012} & \textbf{0.003} & 0.025 & 0.975 & 0.005 \\
\bottomrule
\end{tabular}
\\[6pt]
\begin{tabular}{@{}lrrrrrrrrr@{}}
\toprule
Step & Cov. & MAE & Median & P95 & Gross & Miss & False & Verdict & F.pass \\
\midrule
0\quad Baseline & 1.00 & 1.597 & 0.138 & 6.117 & 0.229 & 0.103 & 0.113 & 0.927 & 0.027 \\
1a\quad DP line only & 1.00 & 0.501 & 0.081 & 0.447 & 0.021 & 0.040 & 0.060 & 0.960 & 0.013 \\
1b\quad Robust depth only & 1.00 & 0.479 & 0.056 & 2.133 & 0.082 & 0.087 & 0.093 & 0.960 & \textbf{0.000} \\
2\quad DP line + robust depth & 1.00 & 0.143 & 0.054 & 0.305 & 0.016 & 0.024 & 0.033 & 0.980 & \textbf{0.000} \\
3\quad + 25 mm cap & 1.00 & 0.143 & 0.054 & 0.305 & 0.016 & 0.024 & 0.033 & 0.980 & \textbf{0.000} \\
4\quad + lower-bound rule & 0.99 & 0.143 & 0.054 & 0.305 & 0.016 & 0.020 & 0.034 & 0.986 & \textbf{0.000} \\
5\quad + risk gate & 0.93 & 0.140 & 0.050 & 0.304 & 0.013 & 0.013 & 0.007 & \textbf{1.000} & \textbf{0.000} \\
6\quad + frame gate \textbf{(recommended)} & 0.91 & 0.141 & \textbf{0.048} & 0.305 & 0.013 & 0.013 & 0.007 & \textbf{1.000} & \textbf{0.000} \\
7\quad Row 2 with on/off line & 1.00 & \textbf{0.131} & 0.055 & \textbf{0.238} & \textbf{0.012} & \textbf{0.000} & \textbf{0.000} & \textbf{1.000} & \textbf{0.000} \\
\bottomrule
\end{tabular}

\end{table}

\begin{figure}[H]
\centering
\includegraphics[width=0.92\linewidth]{\figs/best-stack/error_vs_depth.png}
\caption{Depth error of every scored groove, baseline (blue) against the recommended stack (orange), on
\texttt{stress} and \texttt{shifted}. Errors beyond $\pm12$\,mm are clipped to the edge.}
\label{fig:scatter}
\end{figure}

What the table says, in order of importance:
\begin{enumerate}
\item \textbf{A better line and a better depth stage each remove about two thirds of the mean error} (rows 1a, 1b).
      On mean error the two cannot be ranked against each other (Section~\ref{sec:audit}); DP tracking removes more of
      the gross errors (36\,\% $\to$ 3\,\%) and nearly all missed grooves.
\item \textbf{They add up.} Together they reach 0.33\,mm on \texttt{stress} and 0.14\,mm on other
      units, where each alone was near 0.5\,mm. They fix different errors.
\item \textbf{The risk gate buys safety, not accuracy.} For 9\,\% fewer answers it removes three quarters of the
      false grooves and lifts verdict accuracy to 0.995.
\item \textbf{The frame gate} barely changes these sets (every frame shows a tyre); its job is to refuse photos of
      something else (Table~\ref{tab:real}).
\item \textbf{Two ideas did nothing.} The 25\,mm plausibility cap fired on 6 of 1{,}200 training frames. The
      ``lower-bound'' verdict rule (treat a groove with a hidden floor as ``at least $x$\,mm'') never fired on any
      set. Both are kept as harmless safety nets, but neither earned its place on the evidence.
\item \textbf{On/off subtraction would roughly halve the remaining stress error} (row 7) with no gate at all.
\end{enumerate}

\paragraph{What the recommended stack still gets wrong.} On \texttt{stress}, ten grooves carry 56\,\% of all
remaining error; without them the MAE would be 0.11\,mm. All ten are deep grooves (5--17\,mm), so none changes a
PASS/FAIL verdict. Their causes: the floor of a deep groove lost under a dim laser in sunlight (3), paint or
saturation hiding the floor (2), a glint streak inside the groove (3), and geometry such as $2^\circ$ of rocking (2).
\paragraph{Safety in counts.} Rates hide how few illegal tyres the sets contain, so here are counts. Across
\texttt{ideal}, \texttt{stress} and \texttt{shifted} there are 41 illegal tyres: the baseline passed 9 of them, the
recommended stack none. On \texttt{legal} (64 illegal tyres) the baseline passed none and the stack passed one: a
1.58\,mm groove read as 1.60\,mm. The device needs a
\emph{guard band}: PASS only above, say, 1.8\,mm, and ``borderline'' between 1.4 and 1.8\,mm.

\section{Overall comparison}\label{sec:results}

Tables~\ref{tab:stress}--\ref{tab:real} put the main variant of every family side by side. All numbers are generated
from the scorer's JSON output by \texttt{paper/make\_tables.py}; none is typed by hand.

\begin{table}[H]
\centering\footnotesize\setlength{\tabcolsep}{3.4pt}
\caption{\texttt{stress}: 200 simulated frames in bad conditions. Depth columns in mm. Bold = best in column
(the gated rows answer fewer frames, see Cov.).}\label{tab:stress}
\begin{tabular}{@{}llrrrrrrrrrr@{}}
\toprule
Group & Algorithm & Cov. & MAE & Median & P95 & Gross & Miss & False & Verdict & F.pass & CPU ms \\
\midrule
Reference & Baseline (Steger ridge) & 1.00 & 1.516 & 0.535 & 6.678 & 0.364 & 0.099 & 0.210 & 0.915 & 0.015 & 25.8 \\
\addlinespace[2pt]
Per-column peak & Centre of gravity & 1.00 & 1.061 & 0.085 & 8.525 & 0.096 & 0.069 & 0.160 & 0.930 & 0.005 & 18.9 \\
 & 2-D Gaussian + parabola & 1.00 & 1.089 & 0.089 & 5.317 & 0.073 & 0.101 & 0.185 & 0.960 & \textbf{0.000} & 24.1 \\
\addlinespace[2pt]
Tracking & Viterbi line tracking & 1.00 & 0.544 & 0.086 & 0.591 & 0.033 & 0.018 & 0.115 & 0.955 & \textbf{0.000} & 30.6 \\
\addlinespace[2pt]
Learned line & U-Net, augmented & 1.00 & 0.668 & 0.106 & 2.187 & 0.060 & 0.057 & 0.125 & 0.945 & \textbf{0.000} & 67.3 \\
 & U-Net, no augmentation & 1.00 & 0.236 & 0.091 & 0.467 & 0.021 & 0.009 & 0.030 & \textbf{0.985} & \textbf{0.000} & 63.2 \\
 & Candidate classifier (GBM) & 1.00 & 0.273 & 0.088 & 0.561 & 0.024 & 0.024 & 0.020 & 0.970 & \textbf{0.000} & 32.4 \\
\addlinespace[2pt]
Depth stage & Baseline line + robust depth & 1.00 & 0.464 & 0.093 & 2.225 & 0.121 & 0.113 & 0.145 & 0.940 & 0.010 & 37.1 \\
\addlinespace[2pt]
Gate & Baseline + quality gate & 0.16 & \textbf{0.136} & \textbf{0.069} & 0.634 & \textbf{0.000} & 0.107 & 0.031 & 0.906 & 0.005 & 71.3 \\
\addlinespace[2pt]
Two frames & Laser on/off subtraction & 1.00 & 0.147 & 0.080 & \textbf{0.396} & 0.009 & \textbf{0.003} & \textbf{0.010} & 0.980 & \textbf{0.000} & 22.0 \\
\addlinespace[2pt]
No geometry & End-to-end CNN & 1.00 & 0.483 & 0.247 & 1.963 & 0.102 & 0.006 & 0.015 & 0.980 & 0.015 & 2.3 \\
\bottomrule
\end{tabular}

\end{table}

\begin{table}[H]
\centering\footnotesize\setlength{\tabcolsep}{3.4pt}
\caption{\texttt{ideal}: 150 simulated frames with the holder locked on a clean tyre.}\label{tab:ideal}
\begin{tabular}{@{}llrrrrrrrrrr@{}}
\toprule
Group & Algorithm & Cov. & MAE & Median & P95 & Gross & Miss & False & Verdict & F.pass & CPU ms \\
\midrule
Reference & Baseline (Steger ridge) & 1.00 & 0.147 & 0.054 & 0.474 & 0.008 & 0.090 & \textbf{0.000} & 0.980 & 0.013 & 28.1 \\
\addlinespace[2pt]
Per-column peak & Centre of gravity & 1.00 & 0.062 & 0.039 & 0.191 & \textbf{0.000} & 0.030 & \textbf{0.000} & 0.987 & 0.013 & 20.2 \\
 & 2-D Gaussian + parabola & 1.00 & 0.060 & 0.039 & 0.176 & \textbf{0.000} & 0.041 & \textbf{0.000} & 0.987 & 0.013 & 20.9 \\
\addlinespace[2pt]
Tracking & Viterbi line tracking & 1.00 & 0.064 & 0.041 & 0.186 & \textbf{0.000} & 0.015 & \textbf{0.000} & 0.987 & 0.013 & 27.5 \\
\addlinespace[2pt]
Learned line & U-Net, augmented & 1.00 & 0.061 & 0.040 & 0.173 & \textbf{0.000} & 0.019 & \textbf{0.000} & 0.987 & 0.013 & 63.9 \\
 & U-Net, no augmentation & 1.00 & 0.055 & 0.032 & 0.156 & \textbf{0.000} & 0.011 & \textbf{0.000} & 0.987 & 0.013 & 56.7 \\
 & Candidate classifier (GBM) & 1.00 & 0.154 & 0.079 & 0.527 & \textbf{0.000} & 0.015 & \textbf{0.000} & 0.987 & 0.013 & 32.1 \\
\addlinespace[2pt]
Depth stage & Baseline line + robust depth & 1.00 & 0.081 & \textbf{0.022} & 0.157 & 0.004 & 0.082 & \textbf{0.000} & 0.987 & \textbf{0.000} & 37.0 \\
\addlinespace[2pt]
Gate & Baseline + quality gate & 0.88 & 0.107 & 0.048 & 0.433 & 0.005 & 0.080 & \textbf{0.000} & 0.985 & 0.013 & 72.8 \\
\addlinespace[2pt]
Two frames & Laser on/off subtraction & 1.00 & 0.062 & 0.037 & 0.192 & \textbf{0.000} & 0.015 & \textbf{0.000} & 0.987 & 0.013 & 18.5 \\
\addlinespace[2pt]
No geometry & End-to-end CNN & 1.00 & 0.100 & 0.074 & 0.277 & \textbf{0.000} & \textbf{0.000} & \textbf{0.000} & \textbf{1.000} & \textbf{0.000} & 2.3 \\
\addlinespace[2pt]
Combined & DP line + robust depth & 1.00 & 0.039 & 0.027 & 0.102 & \textbf{0.000} & \textbf{0.000} & \textbf{0.000} & \textbf{1.000} & \textbf{0.000} & 37.7 \\
 & \textbf{Recommended stack} (+ gates) & 1.00 & 0.039 & 0.027 & 0.102 & \textbf{0.000} & \textbf{0.000} & \textbf{0.000} & \textbf{1.000} & \textbf{0.000} & 102.8 \\
 & On/off line + robust depth & 1.00 & \textbf{0.030} & \textbf{0.022} & \textbf{0.085} & \textbf{0.000} & \textbf{0.000} & \textbf{0.000} & \textbf{1.000} & \textbf{0.000} & 28.3 \\
\bottomrule
\end{tabular}

\end{table}

\begin{table}[H]
\centering\footnotesize\setlength{\tabcolsep}{3.4pt}
\caption{\texttt{legal}: 150 worn tyres around the 1.6\,mm limit. Here F.pass is the number that matters.}\label{tab:legal}
\begin{tabular}{@{}llrrrrrrrrrr@{}}
\toprule
Group & Algorithm & Cov. & MAE & Median & P95 & Gross & Miss & False & Verdict & F.pass & CPU ms \\
\midrule
Reference & Baseline (Steger ridge) & 1.00 & 0.053 & 0.030 & 0.210 & \textbf{0.000} & 0.054 & \textbf{0.000} & \textbf{0.993} & \textbf{0.000} & 25.8 \\
\addlinespace[2pt]
Per-column peak & Centre of gravity & 1.00 & 0.024 & 0.019 & 0.064 & \textbf{0.000} & 0.054 & \textbf{0.000} & \textbf{0.993} & \textbf{0.000} & 20.8 \\
 & 2-D Gaussian + parabola & 1.00 & 0.029 & 0.025 & 0.068 & \textbf{0.000} & 0.054 & \textbf{0.000} & \textbf{0.993} & \textbf{0.000} & 22.8 \\
\addlinespace[2pt]
Tracking & Viterbi line tracking & 1.00 & 0.032 & 0.026 & 0.072 & \textbf{0.000} & 0.051 & \textbf{0.000} & \textbf{0.993} & \textbf{0.000} & 27.9 \\
\addlinespace[2pt]
Learned line & U-Net, augmented & 1.00 & 0.036 & 0.033 & 0.087 & \textbf{0.000} & 0.054 & \textbf{0.000} & 0.987 & 0.007 & 62.8 \\
 & U-Net, no augmentation & 1.00 & 0.035 & 0.025 & 0.079 & \textbf{0.000} & 0.054 & \textbf{0.000} & \textbf{0.993} & \textbf{0.000} & 60.1 \\
 & Candidate classifier (GBM) & 1.00 & 0.044 & 0.028 & 0.083 & \textbf{0.000} & 0.051 & \textbf{0.000} & 0.987 & \textbf{0.000} & 25.0 \\
\addlinespace[2pt]
Depth stage & Baseline line + robust depth & 1.00 & \textbf{0.019} & 0.016 & \textbf{0.047} & \textbf{0.000} & \textbf{0.000} & \textbf{0.000} & \textbf{0.993} & \textbf{0.000} & 37.5 \\
\addlinespace[2pt]
Gate & Baseline + quality gate & 0.96 & 0.053 & 0.030 & 0.212 & \textbf{0.000} & 0.035 & \textbf{0.000} & \textbf{0.993} & \textbf{0.000} & 77.4 \\
\addlinespace[2pt]
Two frames & Laser on/off subtraction & 1.00 & 0.026 & 0.020 & 0.066 & \textbf{0.000} & 0.054 & \textbf{0.000} & 0.987 & 0.007 & 22.7 \\
\addlinespace[2pt]
No geometry & End-to-end CNN & 1.00 & 0.072 & 0.062 & 0.187 & \textbf{0.000} & 0.003 & \textbf{0.000} & 0.967 & 0.027 & 2.3 \\
\addlinespace[2pt]
Combined & DP line + robust depth & 1.00 & 0.024 & 0.019 & 0.061 & \textbf{0.000} & \textbf{0.000} & \textbf{0.000} & 0.987 & 0.007 & 39.4 \\
 & \textbf{Recommended stack} (+ gates) & 1.00 & 0.024 & 0.019 & 0.061 & \textbf{0.000} & \textbf{0.000} & \textbf{0.000} & 0.987 & 0.007 & 101.2 \\
 & On/off line + robust depth & 1.00 & \textbf{0.019} & \textbf{0.014} & 0.051 & \textbf{0.000} & \textbf{0.000} & \textbf{0.000} & \textbf{0.993} & \textbf{0.000} & 30.5 \\
\bottomrule
\end{tabular}

\end{table}

\begin{table}[H]
\centering\footnotesize\setlength{\tabcolsep}{3.4pt}
\caption{\texttt{shifted}: three other holder units, each with its own calibration.}\label{tab:shifted}
\begin{tabular}{@{}llrrrrrrrrrr@{}}
\toprule
Group & Algorithm & Cov. & MAE & Median & P95 & Gross & Miss & False & Verdict & F.pass & CPU ms \\
\midrule
Reference & Baseline (Steger ridge) & 1.00 & 1.597 & 0.138 & 6.117 & 0.229 & 0.103 & 0.113 & 0.927 & 0.027 & 24.7 \\
\addlinespace[2pt]
Per-column peak & Centre of gravity & 1.00 & 1.199 & 0.077 & 8.984 & 0.095 & 0.087 & 0.187 & 0.907 & 0.027 & 19.3 \\
 & 2-D Gaussian + parabola & 1.00 & 1.059 & 0.075 & 7.237 & 0.095 & 0.123 & 0.180 & 0.907 & 0.027 & 21.2 \\
\addlinespace[2pt]
Tracking & Viterbi line tracking & 1.00 & 0.501 & 0.081 & 0.447 & 0.021 & 0.040 & 0.060 & 0.960 & 0.013 & 30.5 \\
\addlinespace[2pt]
Learned line & U-Net, augmented & 1.00 & 0.416 & 0.071 & 0.567 & 0.047 & 0.079 & 0.093 & 0.940 & 0.020 & 62.9 \\
 & U-Net, no augmentation & 1.00 & 0.132 & 0.065 & \textbf{0.358} & 0.016 & \textbf{0.028} & 0.047 & 0.973 & 0.013 & 63.0 \\
 & Candidate classifier (GBM) & 1.00 & 0.466 & 0.093 & 1.766 & 0.066 & 0.040 & 0.047 & 0.953 & 0.020 & 29.0 \\
\addlinespace[2pt]
Depth stage & Baseline line + robust depth & 1.00 & 0.479 & \textbf{0.056} & 2.133 & 0.082 & 0.087 & 0.093 & 0.960 & \textbf{0.000} & 35.1 \\
\addlinespace[2pt]
Gate & Baseline + quality gate & 0.52 & \textbf{0.117} & 0.067 & 0.414 & \textbf{0.000} & 0.073 & \textbf{0.000} & \textbf{0.987} & 0.007 & 63.8 \\
\addlinespace[2pt]
Two frames & Laser on/off subtraction & 1.00 & 0.182 & 0.073 & 0.383 & 0.008 & \textbf{0.028} & \textbf{0.000} & 0.987 & 0.013 & 22.3 \\
\addlinespace[2pt]
No geometry & End-to-end CNN & 1.00 & 5.216 & 3.231 & 19.888 & 0.754 & 0.115 & 1.953 & 0.907 & 0.087 & 2.7 \\
\bottomrule
\end{tabular}

\end{table}

\begin{table}[H]
\centering\footnotesize\setlength{\tabcolsep}{3.4pt}
\caption{\texttt{handheld}: 150 simulated frames built to resemble the real photos; errors in pixels of line shift.
The end-to-end network outputs millimetres only and cannot be scored here.}\label{tab:handheld}
\begin{tabular}{@{}llrrrrrrr@{}}
\toprule
Group & Algorithm & Cov. & Line F1 & MAE px & Median px & Gross & Miss & False \\
\midrule
Reference & Baseline (Steger ridge) & 1.00 & 0.385 & 8.11 & 2.81 & 0.494 & 0.466 & 0.327 \\
\addlinespace[2pt]
Per-column peak & Centre of gravity & 1.00 & 0.757 & 2.94 & 0.54 & 0.196 & 0.061 & 0.740 \\
 & 2-D Gaussian + parabola & 1.00 & 0.841 & 2.63 & 0.51 & 0.127 & 0.080 & 0.427 \\
\addlinespace[2pt]
Tracking & Viterbi line tracking & 1.00 & \textbf{0.905} & 3.66 & \textbf{0.40} & 0.103 & \textbf{0.043} & 0.233 \\
\addlinespace[2pt]
Learned line & U-Net, augmented & 1.00 & 0.903 & 1.63 & 0.56 & 0.077 & 0.049 & 0.047 \\
 & U-Net, no augmentation & 1.00 & 0.744 & 1.55 & 0.69 & 0.115 & 0.092 & 0.253 \\
 & Candidate classifier (GBM) & 1.00 & 0.516 & 1.87 & 0.45 & 0.133 & 0.632 & 0.007 \\
\addlinespace[2pt]
Depth stage & Baseline line + robust depth & 1.00 & 0.385 & 7.06 & 0.61 & 0.354 & 0.515 & 0.300 \\
\addlinespace[2pt]
Gate & Baseline + quality gate & 0.34 & 0.785 & 3.77 & 1.96 & 0.294 & 0.056 & 0.196 \\
\addlinespace[2pt]
Two frames & Laser on/off subtraction & 1.00 & 0.753 & \textbf{1.11} & 0.56 & \textbf{0.045} & \textbf{0.043} & 0.453 \\
\addlinespace[2pt]
No geometry & End-to-end CNN & 0.99 & -- & -- & -- & -- & 1.000 & \textbf{0.000} \\
\bottomrule
\end{tabular}

\end{table}

\begin{table}[H]
\centering\footnotesize\setlength{\tabcolsep}{3.4pt}
\caption{Real photos (86). ``$\Delta$\,px'' is the mean difference from the baseline-derived reference on six
truck-tyre frames, so the baseline's own row is near zero by construction. ``Usable kept'': share of the 19 good
photos answered. ``F.meas'': share of unusable photos on which a depth was still reported (55 with no laser,
12 with a laser that cannot be measured). The gate rows are in-sample for the frame gate; cross-validated values
are in the text. Methods needing a laser-off frame or a calibration are omitted.}\label{tab:real}
\begin{tabular}{@{}llrrrrrrr@{}}
\toprule
Group & Algorithm & Line F1 & $\Delta$ px & Found & Within 3 px & Usable kept & F.meas (no laser) & F.meas (laser) \\
\midrule
Reference & Baseline (Steger ridge) & \textbf{1.000} & 0.13 & 6/6 & 6/6 & \textbf{0.947} & 0.691 & 0.750 \\
\addlinespace[2pt]
Per-column peak & Centre of gravity & 0.891 & 5.11 & 5/6 & 3/6 & 0.789 & 0.709 & 0.667 \\
 & 2-D Gaussian + parabola & 0.926 & 1.14 & 6/6 & 5/6 & 0.895 & 0.818 & 0.667 \\
\addlinespace[2pt]
Tracking & Viterbi line tracking & 0.943 & 1.13 & 6/6 & 5/6 & 0.842 & 0.545 & 0.833 \\
\addlinespace[2pt]
Learned line & U-Net, augmented & 0.920 & 1.08 & 6/6 & 5/6 & 0.895 & 0.455 & 0.667 \\
 & U-Net, no augmentation & 0.895 & 7.15 & 6/6 & 3/6 & 0.895 & 0.455 & 0.333 \\
 & Candidate classifier (GBM) & 0.869 & 10.93 & 2/6 & 1/6 & 0.579 & 0.473 & \textbf{0.083} \\
\addlinespace[2pt]
Depth stage & Baseline line + robust depth & \textbf{1.000} & 1.19 & 6/6 & 6/6 & \textbf{0.947} & 0.636 & 0.583 \\
\addlinespace[2pt]
Gate & Baseline + quality gate & \textbf{1.000} & 0.13 & 6/6 & 6/6 & 0.895 & 0.036 & 0.250 \\
\addlinespace[2pt]
Combined & DP line + robust depth & 0.943 & 0.95 & 6/6 & 6/6 & 0.842 & 0.545 & 0.833 \\
 & \textbf{Recommended stack} (+ gates) & 0.943 & 0.95 & 6/6 & 6/6 & 0.789 & \textbf{0.018} & 0.250 \\
\bottomrule
\end{tabular}

\end{table}

\begin{figure}[H]
\centering
\includegraphics[width=0.88\linewidth]{\figs/best-stack/real_overlay.png}
\caption{Baseline (blue) and recommended (orange) lines on three real photos. On the truck tyre they agree to within
a pixel. On the cluttered rover photo both follow background edges; this is the kind of frame the gate must refuse.}
\label{fig:realoverlay}
\end{figure}

\paragraph{Reading the real-photo table.} Three things can be said. First, the recommended stack agrees with the
reference to about 1\,px on all six truck frames, which at the holder's scale would be about 0.5\,mm; whether the
stack or the reference is nearer the truth is unknown. Second, without a gate \emph{every} method reports depths on
half or more of the photos that show no laser at all. Third, with the gate, and with the frame gate re-fitted
without the photos being scored (leave-one-session-out), the stack answers 15 of 19 usable photos, 0 of 55 photos
with no laser, and 3 of 12 photos with an unusable laser.
